% =============================================================================
% Project : THE PHILOSOPHY OF VIABILITY
% Book    : Book I — Ontology of Viability
% Book ID : VS.B01
% Part    : Part I — Identity and Manifestation
% Part ID : VS.B01.P01
% Chapter : Chapter 08 — Comparative Philosophical Examination
% Chapter ID : VS.B01.P01.C08
%
% File :
%   4_ASSETS/diagrams/
%   VS.B01.P01.C08.S16.SS17.P004.DG01.ibn_arabi_and_book_i_relational_distinctions.tex
%
% Asset : DG01
% Structural Identifier :
%   VS.B01.P01.C08.S16.SS17.P004.DG01
%
% Purpose :
%   To represent the principal relations in Ibn Arabi's ontology
%   and distinguish them from the foundational relation between
%   Being and Identity in Book I.
%
% Philosophical Limits :
%   1. The diagram is a conceptual reconstruction, not a literal
%      reproduction of a cosmological diagram by Ibn Arabi.
%   2. Connections do not necessarily represent temporal or
%      efficient-causal sequences.
%   3. Fixed entities are not identified with Book I Identity.
%   4. Divine names are not identified with Determination.
%   5. Tajalli is not identified with generic Manifestation.
%   6. The diagram does not establish an equivalence between
%      the Barzakh and the three mediating hubs of Book I.
%
% Dependencies :
%   tikz, positioning, arrows.meta.
%   (preamble: \usetikzlibrary{positioning, arrows.meta})
%
% =============================================================================

\begin{figure}[htbp]
\centering

\begin{tikzpicture}[
    font=\small,
    every node/.style={inner sep=7pt},
    realnode/.style={
        draw=ViabilityDarkBlue, fill=ViabilityPaleBlue,
        rounded corners=3pt, line width=1pt,
        text width=10.5cm, align=center, minimum height=1.5cm
    },
    namesnode/.style={
        draw=ViabilityGold, fill=ViabilityPaleGold,
        rounded corners=2pt, line width=0.8pt,
        text width=3.5cm, align=center, minimum height=2.8cm
    },
    entitiesnode/.style={
        draw=ViabilityBlue, fill=ViabilityPaleBlue,
        rounded corners=2pt, line width=0.8pt,
        text width=3.5cm, align=center, minimum height=2.8cm
    },
    receptivitynode/.style={
        draw=ViabilityGreen, fill=ViabilityPaleGreen,
        rounded corners=2pt, line width=0.8pt,
        text width=3.5cm, align=center, minimum height=2.8cm
    },
    manifestationnode/.style={
        draw=ViabilityDarkBlue, fill=ViabilityPaper,
        rounded corners=3pt, line width=1pt,
        text width=10.5cm, align=center, minimum height=1.6cm
    },
    booknode/.style={
        draw=ViabilityDarkBlue, fill=ViabilityPaleBlue,
        rounded corners=2pt, line width=0.8pt,
        text width=3.2cm, align=center, minimum height=1.2cm
    },
    identitynode/.style={
        draw=ViabilityGold, fill=ViabilityPaleGold,
        rounded corners=2pt, line width=0.8pt,
        text width=3.2cm, align=center, minimum height=1.2cm
    },
    relation/.style={
        -{Stealth[length=2mm]},
        draw=ViabilityDarkBlue, line width=0.75pt
    },
    conceptual/.style={
        densely dashed, draw=ViabilityGray, line width=0.8pt
    },
    title/.style={
        font=\small\bfseries, align=center, text width=12cm
    },
    label/.style={
        font=\scriptsize, inner sep=2pt, fill=white
    },
    note/.style={
        draw=ViabilityLightGray, fill=white, rounded corners=2pt,
        text width=11cm, align=center
    }
]

% ---------------------------------------------------------------------
% PANEL A — IBN ARABI
% ---------------------------------------------------------------------

\node[realnode] (real) at (0,0) {
    \textbf{The Real — al-Haqq}\\[3pt]
    Non-delimited Existence (\textit{al-wujud al-mutlaq})
};

\node[title, above=2mm of real] (titleA)
    {Ibn Arabi: The Real and Its Determinate Manifestations};

\node[entitiesnode, below=1.6cm of real] (entities) {
    \textbf{Fixed Entities}\\[4pt]
    \textit{A'yan thabita}\\
    Determinate realities\\
    Known in divine knowledge
};

\node[namesnode, left=6mm of entities] (names) {
    \textbf{Divine Names}\\[4pt]
    Relational meanings\\
    Divine attributes\\
    Distinct aspects of disclosure
};

\node[receptivitynode, right=6mm of entities] (receptivity) {
    \textbf{Preparedness and Receptivity}\\[4pt]
    \textit{Isti'dad / Qabiliyya}\\
    Specific capacity of the receiving loci
};

\node[manifestationnode, below=1.4cm of entities] (tajalli) {
    \textbf{Tajalli — Self-Disclosure}\\[4pt]
    The manifestation of the Real under determinate relations
    and according to the receptivity of the loci
};

\node[manifestationnode, below=1.2cm of tajalli] (forms) {
    \textbf{Delimited Manifestations}\\[4pt]
    Existent entities (\textit{a'yan mawjuda})\\
    Differentiated forms of manifested existence
};

% Relations from the Real (vertical drops)

\draw[relation] (names.north |- real.south) -- (names.north)
    node[pos=0.5, left, label, align=right] {relational\\meanings};

\draw[relation] (entities.north |- real.south) -- (entities.north)
    node[pos=0.5, right, label] {divine knowledge};

\draw[relation] (receptivity.north |- real.south) -- (receptivity.north)
    node[pos=0.5, right, label, align=left] {relation to\\the loci};

% Relations to self-disclosure (dashed, non-causal)

\draw[conceptual] (names.south) -- (names.south |- tajalli.north);
\draw[conceptual] (entities.south) -- (tajalli.north);
\draw[conceptual] (receptivity.south) -- (receptivity.south |- tajalli.north);

\draw[relation] (tajalli) -- (forms)
    node[pos=0.5, right, label] {manifested forms};

% ---------------------------------------------------------------------
% PANEL B — BOOK I
% ---------------------------------------------------------------------

\node[title, below=1.5cm of forms] (titleB)
    {Book I: The Foundational Relation};

\node[booknode, below=5mm of titleB, xshift=-4.2cm] (being)
    {\textbf{Being}};

\node[identitynode, below=5mm of titleB, xshift=4.2cm] (identity)
    {\textbf{Identity}};

\draw[conceptual] (being.east) -- (identity.west)
    node[pos=0.5, label, text width=3.3cm, align=center] {
        Analytical Distinction\\
        Ontological Inseparability\\
        Absence of Ontological Priority
    };

\coordinate (axisB) at (titleB |- being.south);

\node[note, below=6mm of axisB] {
    \textbf{Non-equivalence}\\[3pt]
    Fixed Entity $\neq$ Identity \quad Divine Name $\neq$ Determination\\
    Tajalli $\neq$ Generic Manifestation \quad
    Barzakh $\neq$ Book I Mediating Hubs
};

\end{tikzpicture}

\caption[
    Ibn Arabi's ontology and the foundational relation
    between Being and Identity in Book I
]{
    Schematic reconstruction of the relations among
    the Real, divine names, fixed entities, receptivity,
    and self-disclosure in Ibn Arabi's ontology.
    The lower panel separately presents the foundational
    relation between Being and Identity in Book I.
    The diagram is comparative rather than genealogical
    or causal; it does not identify the corresponding
    concepts across the two systems. The connections
    indicate conceptual relations within the reconstructed
    account and do not assert a single temporal or
    efficient-causal sequence.
}

\label{fig:VS.B01.P01.C08.S16.SS17.P004.DG01}

\end{figure}

% =============================================================================
% END OF DIAGRAM DG01
% =============================================================================